\SetOlympString{problem}{%
    \ifOlympLangRu Задача\fi%
    \ifOlympLangKz Есеп\fi%
}
\SetOlympString{solutions}{%
    \ifOlympLangRu Решения\fi%
    \ifOlympLangKz Шешiмдер\fi%
}
\SetOlympString{theory}{Теория}
\SetOlympString{experiment}{Эксперимент}

% ------------------------------------------------------------------
% logo / icon_theor / icon_exp -- optional \includegraphics paths,
% identical in both languages (no \ifOlympLangRu/\ifOlympLangKz
% branching needed), so title_*.tex can \includegraphics{\OlympString{...}}
% instead of hard-coding a path. Relative to wherever the root document
% actually compiles from (<family>/<stage>/<year>/<grade>/<language>/,
% always 5 path segments below shared/example/ for THIS family) -- not
% an absolute-from-repo-root path like "competitions/..." or
% "shared/example/...", which only works in THIS repo's own layout.
% Someone who takes just shared/ into their own project keeps this
% working as long as they preserve the family/stage/year/grade/language
% nesting depth; the file names themselves can freely change year to
% year without ever touching this line, since the relative distance
% to them never does.
%
% logo lives one level above the family folder (shared across every
% stage/year of respa, e.g. an organizer logo); icon_theor/icon_exp
% live at the YEAR folder (this year's own artwork, replaced yearly at
% the same relative path, never renamed).
%
% Not every competition/style HAS a logo or tour icon -- title_*.tex
% uses \OlympIfStringDefined so leaving one of these commented out (or
% never defining it) skips that \includegraphics entirely instead of
% erroring on a missing file.

% ====================================================================

% IF YOU HAVE ANY LOGO/ICON, INCLUDE THE RELATIVE PATH THERE

\SetOlympString{logo}{../../../../../daryn_logo.png}
\SetOlympString{icon_theor}{../../icon_theor.png}
\SetOlympString{icon_exp}{../../icon_exp.png}

% Элемент "../" указывает на переход назад в иерархии папок

% ====================================================================

\SetOlympString{answersheets}{%
\ifOlympLangRu Листы ответов\fi%
\ifOlympLangKz Жауап парақтары\fi%
}

\SetOlympString{part}{%
\ifOlympLangRu Часть\fi%
\ifOlympLangKz Бөлім\fi%
}

\SetOlympString{answersheetsheader}{%
\ifOlympLangRu%
    \ifOlympTheory Листы ответов Теоретического тура\fi%
    \ifOlympExperiment Листы ответов Экспериментального тура\fi%
\fi%
\ifOlympLangKz%
    \ifOlympTheory Теориялық турдың Жауап парақтары\fi%
    \ifOlympExperiment Эксперименталдық турдың Жауап парақтары\fi%
\fi%
}

\SetOlympString{answersheetheadersubtitle}{%
\ifOlympLangRu Не забудьте сдать Ваши листы ответов вместе с чистовиками!\fi%
\ifOlympLangKz Жауап парақтарын таза парақтармен бірге тапсыруды ұмытпаңыз!\fi%
}

\SetOlympString{solutionsheader}{%
    \ifOlympLangRu%
        \ifOlympTheory РЕШЕНИЯ ЗАДАЧ ТЕОРЕТИЧЕСКОГО ТУРА\fi%
        \ifOlympExperiment РЕШЕНИЕ ЗАДАЧИ ЭКСПЕРИМЕНТАЛЬНОГО ТУРА\fi%
    \fi%
    \ifOlympLangKz%
        \ifOlympTheory ТЕОРИЯЛЫҚ ТУРДЫҢ ШЕШIМI\fi%
        \ifOlympExperiment ЭКСПЕРИМЕНТАЛДЫҚ ТУРДЫҢ ШЕШIМI\fi%
    \fi%
}
\SetOlympString{figlabel}{%
    \ifOlympLangRu%
        \ifOlympProblems Рис.~\theproblemnum\printcounter{subproblemnum}.\thefigcounter.\fi%
        \ifOlympSolutions Рис.~\thesolnum\printcounter{subsolnum}.\thefigcounter.\fi%
        \ifOlympMarking Рис.~\thesolnum\printcounter{subsolnum}.\thefigcounter.\fi%
        \ifOlympJury Рис.~\thesolnum\printcounter{subsolnum}.\thefigcounter.\fi%
    \fi%
    \ifOlympLangKz%
        \ifOlympProblems \theproblemnum\printcounter{subproblemnum}.\thefigcounter-сурет.\fi%
        \ifOlympSolutions \thesolnum\printcounter{subsolnum}.\thefigcounter-сурет.\fi%
        \ifOlympMarking \thesolnum\printcounter{subsolnum}.\thefigcounter-сурет.\fi%
        \ifOlympJury \thesolnum\printcounter{subsolnum}.\thefigcounter-сурет.\fi%
    \fi%
}

\SetOlympString{tourname}{%
    \ifOlympLangRu%
        \ifOlympTheory Теоретический тур\fi%
        \ifOlympExperiment Экспериментальный тур\fi%
    \fi%
    \ifOlympLangKz%
        \ifOlympTheory Теориялық тур\fi%
        \ifOlympExperiment Эксперименталдық тур\fi%
    \fi%
}
\SetOlympString{tournamelowercase}{%
    \ifOlympLangRu%
        \ifOlympTheory теоретический тур\fi%
        \ifOlympExperiment экспериментальный тур\fi%
    \fi%
    \ifOlympLangKz%
        \ifOlympTheory теориялық тур\fi%
        \ifOlympExperiment эксперименталдық тур\fi%
    \fi%
}
\SetOlympString{page_abbrev}{%
    \ifOlympLangRu с.\fi%
    \ifOlympLangKz б.\fi%
}
\SetOlympString{organizer}{%
    \ifOlympLangRu РЕСПУБЛИКАНСКИЙ НАУЧНО-ПРАКТИЧЕСКИЙ ЦЕНТР <<ДАРЫН>>\fi%    % competitions.yaml: organizer.uppercase.ru
    \ifOlympLangKz <<ДАРЫН>> РЕСПУБЛИКАЛЫҚ ҒЫЛЫМИ-ПРАКТИКАЛЫҚ ОРТАЛЫҒЫ\fi%  % competitions.yaml: organizer.uppercase.kz
}
\SetOlympString{firstpagesubtitle}{%
    \ifOlympLangRu \OlympString{name}. \OlympString{tourname}, \OlympGrade{} \OlympString{grade_text}. г.~\OlympString{city}, \OlympYear{}\fi%
    \ifOlympLangKz \OlympString{name}. \OlympString{tourname}, \OlympGrade{} \OlympString{grade_text}. \OlympString{city} қ-сы, \OlympYear{}\fi%
}

\SetOlympString{soljanka}{%
    \ifOlympLangRu Солянка\fi%
    \ifOlympLangKz Қоспа\fi%
}

\SetOlympString{soljanka_header}{%
    \ifOlympLangRu Данная задача состоит из трёх независимых частей.\fi%
    \ifOlympLangKz Бұл есеп үш бөлек бөлімнен тұрады.\fi%
}
\SetOlympString{points_short}{%
    \ifOlympLangRu балла\fi%
    \ifOlympLangKz ұпай\fi%
}
\SetOlympString{points_long}{%
    \ifOlympLangRu баллов\fi%
    \ifOlympLangKz ұпай\fi%
}
\SetOlympString{points_answersheet}{%
    \ifOlympLangRu б.\fi%
    \ifOlympLangKz ұ.\fi%
}

\SetOlympString{marking_content}{%
    \ifOlympLangRu \textbf{Содержание}\fi%
    \ifOlympLangKz \textbf{Мазмұны}\fi%
}
\SetOlympString{marking_total}{%
    \ifOlympLangRu \textbf{Итого}\fi%
    \ifOlympLangKz \textbf{Барлығы}\fi%
}
\SetOlympString{eq_formula}{Формула}   % identical in both languages
\SetOlympString{eq_numerical}{%
    \ifOlympLangRu Численное значение\fi%
    \ifOlympLangKz Сандық мән\fi%
}
\SetOlympString{marking_points}{%
    \ifOlympLangRu \textbf{Балл}\fi%
    \ifOlympLangKz \textbf{Ұпайлар}\fi%
}
\SetOlympString{marking_scheme}{Марк-схема}   % identical in both languages
\SetOlympString{marking_schemes}{%
    \ifOlympLangRu Марк-схемы\fi%
    \ifOlympLangKz Марк-схемалар\fi%
}   % identical in both languages
\SetOlympString{grade_text}{%
    \ifOlympLangRu класс\fi%
    \ifOlympLangKz сынып\fi%
}
\SetOlympString{marking_jury}{%
    \ifOlympLangRu Проверяющий\fi%
    \ifOlympLangKz Тексеруші\fi%
}
\SetOlympString{marking_cipher}{Шифр}   % identical in both languages
\SetOlympString{continued}{%
    \ifOlympLangRu (продолжение)\fi%
    \ifOlympLangKz (жалғастыру)\fi%
}
\SetOlympString{writenumbers}
{%
    \ifOlympLangRu Запишите цифры от 0 до 9 в следующую таблицу:\fi%
    \ifOlympLangKz 0-ден 9-ға дейінгі сандарды келесі кестеге жазыңыз:\fi%
}